\documentclass[a4paper,11pt,twoside]{ThesisStyle}

\usepackage{amsmath,amssymb}             % AMS Math
% \usepackage[french]{babel}
\usepackage[latin1]{inputenc}
\usepackage[T1]{fontenc}
\usepackage[left=1.5in,right=1.3in,top=1.1in,bottom=1.1in,includefoot,includehead,headheight=13.6pt]{geometry}
\usepackage{enumerate}
\usepackage{color}
\renewcommand{\baselinestretch}{1.05}

% Table of contents for each chapter

\usepackage[nottoc, notlof, notlot]{tocbibind}
\usepackage{minitoc}
\setcounter{minitocdepth}{2}
\mtcindent=15pt
% Use \minitoc where to put a table of contents

\usepackage{aecompl}

% Glossary / list of abbreviations

\usepackage[intoc]{nomencl}
\renewcommand{\nomname}{List of Abbreviations}

\makenomenclature

% My pdf code

\usepackage{ifpdf}

\ifpdf
  \usepackage[pdftex]{graphicx}
  \DeclareGraphicsExtensions{.jpg}
  \usepackage[a4paper,pagebackref,hyperindex=true]{hyperref}
\else
  \usepackage{graphicx}
  \DeclareGraphicsExtensions{.ps,.eps}
  \usepackage[a4paper,dvipdfm,pagebackref,hyperindex=true]{hyperref}
\fi

\graphicspath{{.}{images/}}

% nicer backref links
\renewcommand*{\backref}[1]{}
\renewcommand*{\backrefalt}[4]{%
\ifcase #1 %
(Not cited.)%
\or
(Cited on page~#2.)%
\else
(Cited on pages~#2.)%
\fi}
\renewcommand*{\backrefsep}{, }
\renewcommand*{\backreftwosep}{ and~}
\renewcommand*{\backreflastsep}{ and~}

% Links in pdf
\usepackage{color}
\definecolor{linkcol}{rgb}{0,0,0.4} 
\definecolor{citecol}{rgb}{0.5,0,0} 

% Change this to change the informations included in the pdf file

% See hyperref documentation for information on those parameters

\hypersetup
{
bookmarksopen=true,
pdftitle="Design and Use of Anatomical Atlases for Radiotherapy",
pdfauthor="Olivier COMMOWICK", 
pdfsubject="Creation of atlases and atlas based segmentation", %subject of the document
%pdftoolbar=false, % toolbar hidden
pdfmenubar=true, %menubar shown
pdfhighlight=/O, %effect of clicking on a link
colorlinks=true, %couleurs sur les liens hypertextes
pdfpagemode=None, %aucun mode de page
pdfpagelayout=SinglePage, %ouverture en simple page
pdffitwindow=true, %pages ouvertes entierement dans toute la fenetre
linkcolor=linkcol, %couleur des liens hypertextes internes
citecolor=citecol, %couleur des liens pour les citations
urlcolor=linkcol %couleur des liens pour les url
}

% definitions.
% -------------------

\setcounter{secnumdepth}{3}
\setcounter{tocdepth}{2}

% Some useful commands and shortcut for maths:  partial derivative and stuff

\newcommand{\pd}[2]{\frac{\partial #1}{\partial #2}}
\def\abs{\operatorname{abs}}
\def\argmax{\operatornamewithlimits{arg\,max}}
\def\argmin{\operatornamewithlimits{arg\,min}}
\def\diag{\operatorname{Diag}}
\newcommand{\eqRef}[1]{(\ref{#1})}

\usepackage{rotating}                    % Sideways of figures & tables
%\usepackage{bibunits}
%\usepackage[sectionbib]{chapterbib}          % Cross-reference package (Natural BiB)
%\usepackage{natbib}                  % Put References at the end of each chapter
                                         % Do not put 'sectionbib' option here.
                                         % Sectionbib option in 'natbib' will do.
\usepackage{fancyhdr}                    % Fancy Header and Footer

% \usepackage{txfonts}                     % Public Times New Roman text & math font
  
%%% Fancy Header %%%%%%%%%%%%%%%%%%%%%%%%%%%%%%%%%%%%%%%%%%%%%%%%%%%%%%%%%%%%%%%%%%
% Fancy Header Style Options

\pagestyle{fancy}                       % Sets fancy header and footer
\fancyfoot{}                            % Delete current footer settings

%\renewcommand{\chaptermark}[1]{         % Lower Case Chapter marker style
%  \markboth{\chaptername\ \thechapter.\ #1}}{}} %

%\renewcommand{\sectionmark}[1]{         % Lower case Section marker style
%  \markright{\thesection.\ #1}}         %

\fancyhead[LE,RO]{\bfseries\thepage}    % Page number (boldface) in left on even
% pages and right on odd pages
\fancyhead[RE]{\bfseries\nouppercase{\leftmark}}      % Chapter in the right on even pages
\fancyhead[LO]{\bfseries\nouppercase{\rightmark}}     % Section in the left on odd pages

\let\headruleORIG\headrule
\renewcommand{\headrule}{\color{black} \headruleORIG}
\renewcommand{\headrulewidth}{1.0pt}
\usepackage{colortbl}
\arrayrulecolor{black}

\fancypagestyle{plain}{
  \fancyhead{}
  \fancyfoot{}
  \renewcommand{\headrulewidth}{0pt}
}

\usepackage{algorithm}
\usepackage[noend]{algorithmic}

%%% Clear Header %%%%%%%%%%%%%%%%%%%%%%%%%%%%%%%%%%%%%%%%%%%%%%%%%%%%%%%%%%%%%%%%%%
% Clear Header Style on the Last Empty Odd pages
\makeatletter

\def\cleardoublepage{\clearpage\if@twoside \ifodd\c@page\else%
  \hbox{}%
  \thispagestyle{empty}%              % Empty header styles
  \newpage%
  \if@twocolumn\hbox{}\newpage\fi\fi\fi}

\makeatother
 
%%%%%%%%%%%%%%%%%%%%%%%%%%%%%%%%%%%%%%%%%%%%%%%%%%%%%%%%%%%%%%%%%%%%%%%%%%%%%%% 
% Prints your review date and 'Draft Version' (From Josullvn, CS, CMU)
\newcommand{\reviewtimetoday}[2]{\special{!userdict begin
    /bop-hook{gsave 20 710 translate 45 rotate 0.8 setgray
      /Times-Roman findfont 12 scalefont setfont 0 0   moveto (#1) show
      0 -12 moveto (#2) show grestore}def end}}
% You can turn on or off this option.
% \reviewtimetoday{\today}{Draft Version}
%%%%%%%%%%%%%%%%%%%%%%%%%%%%%%%%%%%%%%%%%%%%%%%%%%%%%%%%%%%%%%%%%%%%%%%%%%%%%%% 

\newenvironment{maxime}[1]
{
\vspace*{0cm}
\hfill
\begin{minipage}{0.5\textwidth}%
%\rule[0.5ex]{\textwidth}{0.1mm}\\%
\hrulefill $\:$ {\bf #1}\\
%\vspace*{-0.25cm}
\it 
}%
{%

\hrulefill
\vspace*{0.5cm}%
\end{minipage}
}

\let\minitocORIG\minitoc
\renewcommand{\minitoc}{\minitocORIG \vspace{1.5em}}

\usepackage{multirow}
\usepackage{slashbox}

\newenvironment{bulletList}%
{ \begin{list}%
	{$\bullet$}%
	{\setlength{\labelwidth}{25pt}%
	 \setlength{\leftmargin}{30pt}%
	 \setlength{\itemsep}{\parsep}}}%
{ \end{list} }

\newtheorem{definition}{Definition}
\renewcommand{\epsilon}{\varepsilon}

% centered page environment

\newenvironment{vcenterpage}
{\newpage\vspace*{\fill}\thispagestyle{empty}\renewcommand{\headrulewidth}{0pt}}
{\vspace*{\fill}}



\begin{document}

\begin{titlepage}
\begin{center}
\noindent {\large \textbf{UNIVERSITY OF NICE - SOPHIA ANTIPOLIS}} \\
\vspace*{0.3cm}
\noindent {\LARGE \textbf{DOCTORAL SCHOOL STIC}} \\
\noindent \textbf{SCIENCES ET TECHNOLOGIES DE L'INFORMATION \\ ET DE LA COMMUNICATION} \\
\vspace*{0.5cm}
\noindent \Huge \textbf{P H D\ \ T H E S I S} \\
\vspace*{0.3cm}
\noindent \large {to obtain the title of} \\
\vspace*{0.3cm}
\noindent \LARGE \textbf{PhD of Science} \\
\vspace*{0.3cm}
\noindent \Large of the University of Nice - Sophia Antipolis \\
\noindent \Large \textbf{Specialty : \textsc{Computer Science}}\\
\vspace*{0.4cm}
\noindent \large {Defended by\\}
\noindent \LARGE Olivier \textsc{Commowick} \\
\vspace*{0.8cm}
\noindent {\Huge \textbf{Design and Use of Aatomical Atlases for Radiotherapy}} \\
\vspace*{0.8cm}
\noindent \Large Thesis Advisor: Gr�goire \textsc{Malandain} \\
\vspace*{0.2cm}
\noindent \Large prepared at INRIA Sophia Antipolis, \textsc{Asclepios} Team\\
\vspace*{0.2cm}
\noindent \large defended on February 14, 2007 \\
\vspace*{0.5cm}
\end{center}
\noindent \large \textbf{Jury :} \\
\begin{center}
\noindent \large 
\begin{tabular}{llcl}
      \textit{Reviewers :}	& Patrick \textsc{Clarysse}		& - & CNRS (CREATIS)\\
				& Louis \textsc{Collins}		& - & McGill University\\
      \textit{Advisor :}	& Gr�goire \textsc{Malandain}		& - & INRIA (Asclepios)\\
      \textit{President :}	& Nicholas \textsc{Ayache}		& - & INRIA (Asclepios)\\
      \textit{Examinators :}   & Pierre-Yves \textsc{Bondiau}          & - & Centre Antoine Lacassagne (Nice)\\
      				& Guido \textsc{Gerig}			& - & University of North Carolina\\
      				& Vincent \textsc{Gr�goire}		& - & Universit� Catholique de Louvain\\
      \textit{Invited :}		& Hanna \textsc{Kafrouni}		& - & DOSISoft S.A.
\end{tabular}
\end{center}
\end{titlepage}
\sloppy

\titlepage


\dominitoc

\pagenumbering{roman}

 \cleardoublepage

\section*{Acknowledgments}

Last thing to do :-)

\tableofcontents

\mainmatter

\chapter{Introduction}
\label{chap:intro}
\minitoc

\section{Illustration Example}

\subsection{A subsection just for fun}

Sorry I won't write your PhD here ;) This small text just to mention that this style supports writing with accents such as in french words (th�se, d�finir, ...). Also I put here a simple way to include an image. This is standard latex. For pdflatex compilation, the extension of the images is jpg. For latex compilation, this is ps or eps. The base folder containing images is set in formatAndDefs.tex, as well as the default extensions added to the image names.

\begin{figure}[!htbp]
  \begin{center}
    \includegraphics[width=0.9\textwidth]{Chapitre1/arctic_control}
  \end{center}
  \caption{A nice image...}
  \label{fig:jolieImage}
\end{figure}

\section{An equation}

Just to show argmin and partial derivative commands.

\begin{equation}
  T = \argmin_T E(T,R,F)
\end{equation}

Regularization:

\begin{equation}
  \pd{T}{t} = \Delta T
\end{equation}

\section{An other section}

Showing a great bullet list environment:

\begin{bulletList}
 \item First point
 \item Second point
% \item Here is an abbreviation reference \nomenclature{DTI}{Diffusion Tensor Imaging} DTI
\end{bulletList}


\chapter{No-name yet}

\newpage

\section{Query taxonomy and re-usability formalization}

This section presents the description and formalization of the different types of queries which \textit{(i)} can be processed by our integration approach; and \textit{(ii)} can be compared to previous integration requests in order to take advantage from previous integration plans. The query definition is introduced below.

\begin{definition}
A query is defined as a $n$-tuple:
%
\begin{center}
$Q := \langle A, R, S, C, w  \rangle$
\end{center}
%
where: $A$ is a set of abstract services defining the query $Q$;
$R$ is a set of user preferences that can be defined over the data services or the entire query;
$S$ is a set of data services that were selected satisfying the restrictions defined by $R$ to potentially rewrite the query $Q$;
$C$ is a set of compositions that were produced using the data services in $S$ and satisfying the restrictions defined by $R$ that potentially can answer the query $Q$; and
$w$ is the composition that were selected and executed to answer the query $Q$. 
\end{definition}

The query taxonomy proposed below is defined according to the type of relation that can be established between two queries. 
Queries are classified in four groups: 
\begin{itemize}
\item \textit{Group 1:} The data denoted by the answer of $Q_{1}$ is the same data expected by the answer of $Q_{2}$. For example, $Q_{1}$ and $Q_{2}$ retrieve patients that were infected by pneumonia.
\item \textit{Group 2:} The data denoted by the answer of $Q_{1}$ is a subset of the data denoted by the answer of $Q_{2}$. For example, $Q_{2}$ retrieves patients that were infected by pneumonia and $Q_{1}$ retrieves patients that were infected by pneumonia and treated by the doctor Lucas.
\item \textit{Group 3:} The data denoted by the answer of $Q_{1}$ is a superset of the data denoted by the answer of $Q_{2}$. For example, $Q_{2}$ retrieves patients that were infected by pneumonia and treated by the doctor Lucas, and $Q_{1}$ retrieves patients that were infected by pneumonia.
\item \textit{Group 4:} The data denoted by the answer of $Q_{1}$ is different of the data denoted by the answer of $Q_{2}$. For example, $Q_{2}$ retrieves patients that were infected by pneumonia and treated by the doctor Lucas, and $Q_{1}$ retrieves patients that were infected by pneumonia with admission in the hospital Edouard Herriot.
\end{itemize}

To understand the different types of query, basic concepts regarding \textit{(i)} user requirements, \textit{(ii)} requirements domain, \textit{(iii)} requirements evaluation and \textit{(iv)} comparable requirements should be introduced:

\begin{definition}
An user requirement $r$ is in the form $x \otimes c$, where $x$ is an identifier; $c$ is a constant; and $\otimes \in\lbrace \geq, \leq, =, \neq, <, >\rbrace$. 
%
The user requirement $r$ could concern \textit{(i)} the entire \textsl{query}, in this case noted as $r_{Q}$; or \textit{(ii)} a single service, noted as $r_{S}$. For instance, the total response time is obtained by adding the response time of each service involved in the composition.
\end{definition}

\begin{definition}
A requirement domain is a set of possible values which can be assumed by an user requirement $r$, represented by $Dom(r)$. 
For instance, a requirement domain ``response time'' includes the possible values associated to the response time user requirement. 
Each user requirement $r_{i}$ has its own requirement domain $D_{i}$. 
\end{definition}

\begin{definition}
The evaluation of an user requirement $r$, indicated by $eval(r)$, returns a set of values $\lbrace v_{1},..,v_{i} \rbrace$ that can be assigned to $r$ such that $\lbrace v_{1},..,v_{i} \rbrace \subset Dom(r)$.
\end{definition}

\begin{definition}
Given two user requirements $r_{1}$ and $r_{2}$, both can be comparable, denoted by $r_{1} \perp r_{2}$, if and only if:  $Dom(r_{1}) = Dom(r_{2})$.
\end{definition}

The thirteen types of queries included in the taxonomy described in the following sections are organized according to their groups.

\subsection{Queries belonging to the Group 1}

Group 1 contains queries of four different types: \textit{(i)} equivalent; \textit{(ii)} subset; \textit{(iii)} superset; and \textit{(iv)} a special case. 

\subsubsection{Query type 1: $Q_{1}$ is equivalent to $Q_{2}$}

The \textit{first} type deals with \textit{equivalent queries}. Two queries $Q_{1}$ and $Q_{2}$ are equivalent when:
%
\begin{enumerate}[a)]
\item They expect the same data as answer, which means they cover the same abstract services. For instance, the set of abstract service of $Q_{1}$, denoted as $Q_{1}.A$, is equals to the set of abstract services of $Q_{2}$, denoted as $Q_{2}.A$.
%
\begin{center}
$Q_{1}.A = Q_{2}.A$
\end{center}
%
\item For each user requirement $r_{i}$ in $Q_{1}.R$, there is a user requirement $r_{j}$ in $Q_{2}.R$ such that the evaluation of $r_{i}$ is equal to the evaluation of $r_{j}$. 
Consequently, the score of $Q_{1}.R$ is equals to the score of $Q_{2}.R$. The equivalence between requirements is defined below.
\end{enumerate}

\begin{definition}\label{def:reqeq}
A set of \textsl{user requirements} $R_{1}$ is equivalent to a set of \textsl{user requirements} $R_{2}$, represented by $R_{1} \equiv R_{2}$, if and only if: $\forall r_{i} \in R_{1}, \ \exists r_{j} \in R_{2} \ \vert \ eval (r_{i}) = eval(r_{j}) \ and \ \vert R_{1} \vert = \vert R_{2} \vert$.
\end{definition}

From the re-usability point of view, all data services filtered to the query $Q_{1}$ which are \textit{online} in the moment, denoted $online(Q_{1}.S)$, could be reused in the query $Q_{2}$. 
Moreover, the rewritings (compositions) produced to the query $Q_{1}$ which use the data services in $online(Q_{1}.S)$, denoted as $available(Q_{1}.C)$, could also be used to answer the query $Q_{2}$. The query type 1 definition is presented below.

\begin{definition}\label{def:qt1}
Query Type 1 -- a query $Q_{1}$ is equivalent to a query $Q_{2}$, if and only if: $Q_{1}.A = Q_{2}.A$ and $Q_{1}.R_{1} \equiv Q_{2}.R_{2}$
\end{definition}

\subsubsection{Query type 2: $Q_{2}$ is a subset of $Q_{1}$}

The \textit{second} type deals with \textit{query subsets} due to more restrict user requirements. Given two queries $Q_{1}$ and $Q_{2}$, $Q_{2}$ is a subset of $Q_{1}$ when:
%
\begin{enumerate}[a)]
\item They expect the same data as answer, which means they cover the same abstract services. 
For instance, the set of abstract service of $Q_{1}$, denoted as $Q_{1}.A$, is equals to the set of abstract services of $Q_{2}$, denoted as $Q_{2}.A$.
%
\begin{center}
$Q_{1}.A = Q_{2}.A$
\end{center}
%
\item For all user requirement $r_{i}$ in $Q_{2}.R$, there is at least one $r_{j}$ in $Q_{1}.R$ such that the evaluation of $r_{i}$ is contained in the evaluation of $r_{j}$. 
For all $r_{k}$ in $Q_{2}.R$, there is no $r_{l}$ in $Q_{1}.R$ such that the evaluation of $r_{l}$ is contained in the evaluation of $r_{k}$. 
Consequently, the score of $Q_{1}.R$ is lower than the score of $Q_{2}.R$. The definition of more restrict requirements is presented below.
\end{enumerate}

\begin{definition}\label{def:reqmore}
Given a set of \textsl{user requirements} $R_{1}$ and $R_{2}$, $R_{1}$ is more restrict than $R_{2}$, represented by $R_{1} \ \rhd \ R_{2}$, if and only if: $\forall r_{i} \in R_{1}, \ \exists r_{j} \in R_{2}, \ \nexists r_{k} \in R_{2} \ \vert \ eval (r_{i}) \subset eval(r_{j}) \ and \ eval (r_{k}) \subset eval(r_{i}) \ and \ \vert R_{1} \vert = \vert R_{2} \vert$.
\end{definition}

From the re-usability point of view, a subset of the data services filtered to the query $Q_{1}$ which are \textit{online} in the moment, $online(Q_{1}.S)$, could be reused in the query $Q_{2}$. This fact occurs due to the more restrict requirements imposed by $Q_{2}$.
With respect to the compositions, a subset of the rewritings produced to the query $Q_{1}$ could also be used to answer the query $Q_{2}$. These rewritings should use the data services in $online(Q_{1}.S)$, denoted as $available(Q_{1}.C)$, and respect the more restrict requirements defined in $Q_{2}$. 
The query type 2 definition is presented below.

\begin{definition}\label{def:qt2}
Query Type 2 -- a query $Q_{1}$ is a subset of a query $Q_{2}$, if and only if: $Q_{1}.A = Q_{2}.A$ and $Q_{1}.R_{1} \rhd Q_{2}.R_{2}$
\end{definition}



\appendix

\chapter{Appendix Example}
\label{chap:appendix1}

\section{Appendix Example section}

And I cite myself to show by bibtex style file (two authors) \cite{Commowick_MICCAI_2007}.

This for other bibtex stye file : only one author \cite{Oakes_RStat_1999} and many authors \cite{Guimond_CVIU_2000}.

\bibliographystyle{ThesisStyle}
\bibliography{bibliography}

%\printnomenclature

\cleardoublepage
\begin{vcenterpage}
\noindent\rule[2pt]{\textwidth}{0.5pt}
\begin{center}
{\large\textbf{Design and Use of Numerical Anatomical Atlases for Radiotherapy\\}}
\end{center}
{\large\textbf{Abstract:}}
The main objective of this thesis is to provide radio-oncology specialists with automatic tools for delineating organs at risk of a patient undergoing a radiotherapy treatment of cerebral or head and neck tumors.
\\
To achieve this goal, we use an anatomical atlas, i.e. a representative anatomy associated to a clinical image representing it. The registration of this atlas allows to segment automatically the patient structures and to accelerate this process. Contributions in this method are presented on three axes.
\\
First, we want to obtain a registration method which is as independent as possible w.r.t. the setting of its parameters. This setting, done by the clinician, indeed needs to be minimal while guaranteeing a robust result. We therefore propose registration methods allowing to better control the obtained transformation, using outlier rejection techniques or locally affine transformations.
\\
The second axis is dedicated to the consideration of structures associated with the presence of the tumor. These structures, not present in the atlas, indeed lead to local errors in the atlas-based segmentation. We therefore propose methods to delineate these structures and take them into account in the registration.
\\
Finally, we present the construction of an anatomical atlas of the head and neck region and its evaluation on a database of patients. We show in this part the feasibility of the use of an atlas for this region, as well as a simple method to evaluate the registration methods used to build an atlas.
\\
All this research work has been implemented in a commercial software (Imago from DOSIsoft), allowing us to validate our results in clinical conditions.
\\
{\large\textbf{Keywords:}}
Atlas-based Segmentation, non rigid registration, radiotherapy, atlas creation
\\
\noindent\rule[2pt]{\textwidth}{0.5pt}
\end{vcenterpage}

\end{document}
